\documentclass{jsarticle}
\usepackage{amsmath}
\usepackage{amssymb}
\begin{document}
\title{}
\author{}
\date{}
\maketitle


　　ブレインインターフェースとその応用、機知の仕組み

　ブレインインターフェースは、機械において、音声を文字列に表示するには、
まず、音声を音階の7階層に分解して、その上に、大脳基底核を経由する共振回路の共鳴を使って、唇の動きのデータを電気信号に変えて、この2種類のデータから、音声を文字列に表示する。頭頂葉から大脳基底核を経由する体感触覚によるデータを使って、そのデータをイメージと言語と文字列に変えて、画像処理を使って、動的に映像やページを処理する。例えば、Wordのページをめくったりする。
ASを頭頂葉から大脳基底核を経由するデータを電気信号に変えて、機知の仕組みを担う第32，33野を経由するドーパミンがそれぞれの役割をする共振回路の共鳴を使って、血液のリボ核酸のヘリコバクターがこの共鳴する共振回路に接続して、ASを動かす。

$$ H_n=\sigma(E_n )×\sigma(K_n), e_n=f^{-1}(x)xf(x)$$
$$e_n=\mathrm{ker}f \oplus \mathrm{im}f,H_n^2=|(|ds^2 |)|,=H_n×K_n $$

$$E_n=σ(K_n )×σ(H_n ),π(σ_1,σ_2 )=  {d \over df} F(x) $$
$$ x\Gamma(x) =\pi(\sigma_1, \sigma_2 )-  (8 \pi G({p \over c^3} + 
{V \over S})) $$
$$=2\int  ||\sin2x||^2 d \tau $$
$$ =4 \int ||\sin x  \cos x||^2 d \tau \to 4V = g_{ij}^2 $$

   固有エントロピー = 原子の固有エントロピー - ３次元多様体のエントロピー


   ブレインインターフェースの大脳基底核の固有エントロピー値

\end{document}
